\documentclass[../main.tex]{subfiles}
\begin{document}
%==========================================TIÊU ĐỀ TẬP========================================
\begin{table}[h]
  \begin{adjustwidth}{-1cm}{}
    \begin{tabular}{>{\centering\arraybackslash}p{6cm}|>{\raggedright\arraybackslash}p{12.5cm}}
      \multirow{5}{6cm}
      % Cột bên trái: ÉPISODE và số tập
      {\\[-40pt]
        \flaregothic\fontsize{20pt}{20pt}\selectfont \centering
        {\contour{errorfive}{\textcolor{black}{ÉPISODE}}}\color{black}\\
        \burbank\fontsize{72pt}{72pt}\selectfont
        \includegraphics[height=3.12359cm]{../../../../Image/Book?/number/num35_alt.png}
        \\[18pt]
        % \flaregothic\fontsize{14pt}{14pt}\selectfont
        % \color{epattr}BIENVENUE, \uline{\color{Banpire} Mel} !
        % \flaregothic\fontsize{14pt}{14pt}\selectfont{\contour{errorfive}{\textcolor{white}{BIENVENUE, }}
        %     {\color{black} Naomi}
        %     {\contour{errorfive}{\textcolor{white} {!}}}}\\
        \sen\fontsize{18pt}{18pt}\selectfont{\contour{errorfive}{\textcolor{white}{Mini-histoire}}}
        \color{black}
      }
       &
      % Dòng thứ nhất: tên tập tiếng Nhật
      {\fotskip\fontsize{18pt}{18pt}\selectfont \ruby{光}{ひかり}の\ruby{欠片}{かけら}}
      \\[6pt]
      % Dòng thứ hai: tên tập tiếng Anh
       & {\cafeteria\fontsize{18pt}{18pt}\selectfont Pieces of Light}               \\[4pt]
      % Dòng thứ ba: tên tập tiếng Việt
       & {\vnmsans\fontsize{18pt}{18pt}\selectfont Những mảnh ghép của ánh sáng}    \\[8pt]
      % Dòng thứ tư: nhân vật xuất hiện
       & {\caviar{\fontsize{14pt}{14pt}\selectfont Nhân vật: Toàn bộ dàn nhân vật}}
      % \\[6pt]
      % Dòng cuối cùng: Thời gian diễn ra của tập (thường sẽ là ngày phát hành)
      % & {\continuum\fontsize{16pt}{16pt}\selectfont Ngày phát hành: 11/06/2022}\\[8pt]
    \end{tabular}
  \end{adjustwidth}
\end{table}
%===========================================NỘI DUNG==========================================
\color{black}
\vspace{30pt}

\begin{center}
  \uline{Ngày thứ hai đầu tuần sau.}
  \pov{Hikawa Kaigou}
\end{center}

Ngày thứ hai đầu tuần sau, guồng quay của trường lớp lại bắt đầu. Với việc Aogami đã lấy lại sự bình yên vốn có, chúng tôi lại tiếp tục với vỏ bọc của những học sinh trung học bình thường.

Trước khi đi học, chúng tôi đã thống nhất sẽ để Nao ở nhà. Lý do rất đơn giản: con bé mới chỉ tám tuổi về mặt sinh học, chưa có bất kỳ giấy tờ tùy thân nào, và dĩ nhiên là chẳng có lý do gì để nhét một đứa trẻ tiểu học vào giữa một lớp cấp ba toàn những cô cậu mười lăm, mười sáu tuổi cả. Hơn nữa, với xuất thân đặc biệt của Nao, chúng tôi nghĩ tốt nhất là nên để con bé có thêm thời gian làm quen với các thiết bị trong nhà (mặc dù tôi cũng hơi lo khi để một đứa trẻ tầm tiểu học ở nhà một mình). Kuon-kun đã cẩn thận bật một đĩa phim hoạt hình, dặn dò Nao ngoan ngoãn ngồi xem và cấm tuyệt đối việc mở cửa cho người lạ. Suzu thì chuẩn bị sẵn đồ ăn trưa trong tủ lạnh. Tôi thừa nhận rằng mình đã tự tin rằng bọn tôi đã chu toàn mọi thứ.

Nhưng có vẻ như tôi đã đánh giá quá thấp sự hiếu kỳ của một đứa trẻ bị nhốt hơn 1300 năm.

Vừa bước tới cửa lớp 1-C, tôi đã cảm nhận được một bầu không khí rộn ràng bất thường. Thay vì tiếng nói chuyện phiếm ồn ào hay tiếng lật sách sột soạt quen thuộc, cả lớp dường như đang xúm tụm lại một góc.

``Trời ơi, dễ thương ghê\~{}'' tiếng của Yae vang lên lanh lảnh.

``Nhìn cặp má phúng phính này đi, mình muốn nựng quá!'' giọng của Kanade tiếp nối.

``Tóc tết hai bím kìa, búp bê sống là đây chứ đâu!'' tiếng của Inari xuýt xoa.

Kuon-kun đi ngay cạnh tôi, nhíu mày khó hiểu: ``\vie{Bọn họ đang bàn luận về con thú cưng nào mới được mang vào lớp à?}''

Tò mò, chúng tôi chen qua đám đông học sinh đang xúm xít quanh bàn của Sakura. Và ngay tại giây phút đó, máu trong người ba chúng tôi dường như đông cứng lại.

Ngồi chễm chệ trên chiếc ghế trống, hai chân nhỏ bé đung đưa không chạm đất, tay ôm một hộp sữa dâu\dots\ chính là Nao! Con bé đang mặc nguyên chiếc váy mùa hè mà Suzu mua cho cuối tuần trước, đôi mắt lục bảo chớp chớp nhìn đám đông bằng vẻ điềm tĩnh vô song.

``\textbf{Hả!?}'' Kuon-kun và Suzu đồng thanh hét lên, còn cằm tôi thì rớt thẳng xuống sàn.

Nao vào tới đây bằng cách nào? Sao con bé lại biết đường tới trường? À không, khoan đã, khi còn là một hồn ma bay lơ lửng, con bé có thể tự do theo dõi chúng tôi mọi lúc mọi nơi, nên việc biết đường tới lớp 1-C chẳng có gì khó hiểu cả. Vấn đề là, tại sao con bé lại mò tới đây thay vì ngoan ngoãn ngồi nhà xem phim hoạt hình!?

Nhìn thấy chúng tôi, Nao khẽ mỉm cười, giơ tay lên vẫy vẫy. ``Onii-chan, Suzuki-onee-chan, Kaigou-onee-chan\~{} Các bạn học của anh chị thú vị thật đấy!''

Sự xuất hiện của ba chúng tôi làm đám đông giãn ra một chút, nhưng cũng chính là lúc màn ra mắt ``thảm họa'' nhất lịch sử bắt đầu diễn ra theo cách không một ai trong bộ ba chúng tôi lường trước được.

Đầu tiên là nhóm ba người Shino, Sakura và Shion. Bọn họ đứng ngay sát Nao, nhưng sắc mặt thì tái mét như vừa bị rút cạn sinh lực. Shino và Sakura run rẩy, lùi lại mấy bước, ngón tay chỉ thẳng vào Nao lắp bắp.

``C-C-Cô bé đó\dots\ đôi mắt đó\dots'' Sakura ôm đầu, vỡ òa. ``Đó là con ma ở gốc cây tối hôm đó mà!!''

``\textbf{Oán linh! Oán linh đi học!!}'' Shino hét lên một tiếng rồi chui tọt xuống gầm bàn trốn. Shion thì đứng đơ như tượng đá, mồ hôi lạnh chảy ròng ròng trên trán.

Nhìn bộ dạng như gặp quỷ giữa ban ngày của họ, Hanabi nhíu mày khó hiểu: ``Ba cậu bị sao vậy? Tự nhiên hét toáng lên thế?''

``\textbf{K-Không có gì đâu!!}'' Sakura và Shino (vọng lên từ gầm bàn) đồng thanh hét lớn, vội vàng xua tay lảng tránh, mồ hôi hột tuôn ra như tắm.

``\dots Được rồi?'' cô nàng tóc nâu nhíu mày, vẫn tỏ vẻ nghi ngờ với hai cô nàng. Shino ơi, Sakura ơi, cứ trốn dưới gầm như thế người ta càng nghi ngờ hơn đấy\dots

Ở một góc khác, Azuki, Nanase và Mari - những người cũng từng vào rừng lần ba và chạm trán Nao - há hốc mồm ngạc nhiên tột độ. Nanase dụi mắt liên tục: ``Không thể nào\dots\ Tôi nhớ là con bé đó trong suốt cơ mà? Sao bây giờ lại có cả bóng kìa?'' Mari thì nắm chặt tay Azuki, lẩm bẩm cầu nguyện liên hồi.

Ayame và Wata đứng cạnh nhau. Hội trưởng hội học sinh - người cũng tham gia nhóm Azuki đêm đó - có chung phản ứng sốc cực độ. Nhưng Hội phó thì lại ngơ ngác nhìn cô bạn mình: ``Cậu sao vậy Ayame? Có một đứa trẻ đi lạc thôi mà, làm gì mà nhìn như thấy ma thế?''

``K-Không có gì đâu! Tôi chỉ hơi mệt xíu thôi!'' Ayame vội vàng đánh trống lảng, cười gượng ép gạt đi.

Nhóm của Touka, Saya và Aku (những người từng đi vào rừng lần thứ hai) thì xôn xao bàn tán.

``Này, đôi mắt xanh lục đó\dots\ trông quen lắm đúng không?'' Touka huých khuỷu tay Saya.

``Đừng nói với mình đây là\dots\ cái bóng trắng ở khu nhà cũ nhé?'' Saya nuốt nước bọt.

``Chắc là người giống người thôi\dots\ làm gì có ma quỷ đi học được\dots'' Aku lẩm bẩm, tự trấn an bản thân dù hai chân đang run lẩy bẩy.

Bất ngờ nhất là Uzuki. Cô nàng Nhiếp ảnh hoảng hốt lôi trong cặp ra bức ảnh chụp ở trong rừng mà trước đây từng chụp được cái bóng của Nao. Uzuki run rẩy cầm bức ảnh so sánh với người thật đang ngồi đung đưa chân trước mặt, rồi thốt lên: ``Hàng thật! Là hàng thật 100\%!!''

Giữa khung cảnh hỗn loạn kẻ thì khen dễ thương, người thì sợ mất mật ấy, Yuka - cô nàng có giác quan thứ sáu nhạy bén nhất lớp - bước tới gần Nao. Cô mỉm cười đầy ẩn ý, khẽ vuốt lọn tóc tết của con bé. ``Một linh hồn thật thuần khiết và trong sạch\dots\ Chúc mừng em đã có một cơ thể mới nhé\~{}'' Yuka nói, một câu nói khiến nhóm của Shino càng thêm rùng mình, còn chúng tôi thì muốn xỉu ngang.

Nghe thấy câu nói đó, Akane bước tới, nhướng mày hỏi: ``Yuka, ý bà là sao? Cơ thể mới là sao?''

Bản thân cô nàng bờm Chó cũng từng thấy bóng mờ của hồn ma Nao trong một lần đi khảo sát, nhưng vì lúc đó hình bóng quá mờ nhạt nên cô cũng không dám chắc người đang ngồi đây có phải là cùng một thực thể đó không.

``Fufu\dots\ Không có gì đâu\~{}'' cô nàng bờm Mèo chỉ đưa tay che miệng cười đáp, ánh mắt vẫn không rời khỏi Nao.

Với những người hoàn toàn không biết chuyện ma quỷ gì như Honoka, Kanade, Yuki và Koishin, sự tò mò của họ là thứ cứu vãn bầu không khí.

Cô nàng bờm Cáo vàng nhào tới, hăng hái vuốt lọn tóc của Nao: ``Trời ơi, em gái nhỏ, em dùng dầu gội gì mà tóc mềm thế này? Cho chị nựng má một cái nha!''

Kanade đứng bên cạnh gật gù, vuốt cằm ra vẻ suy tư: ``Chắc chắn là gen gia đình tốt rồi, nhìn thần thái này quả là hiếm có khó tìm.''

Koishin hừ mũi, quay mặt đi chỗ khác dù hai vành tai đã đỏ ửng: ``Hmph! Trẻ con ồn ào chết đi được. Tôi chẳng thấy dễ thương chỗ nào.''

Yuki lập tức huých cùi chỏ vào mạn sườn cậu ta, nhếch mép mỉa mai: ``Khỏi giả bộ đi Koishin. Ai chẳng biết lúc nãy cậu vừa lén lấy điện thoại ra tính chụp trộm người ta à\~{}''

``C-Cái con nhỏ này! Ai thèm chụp trộm!?'' cô nàng Tsundere giãy nảy lên cãi lại.

Hotaru bước lên, chống tay xuống bàn, thân thiện hỏi: ``Này nhóc, em từ đâu tới vậy? Bị lạc người thân à?''

Tim tôi, Kuon-kun và Suzu lập tức nhảy lên tận họng. Chúng tôi lao lên định bịt miệng Nao lại để tìm đại một cái cớ nào đó, nhưng con bé đã nhanh miệng cất lời trước, chất giọng trầm tĩnh và vô cùng tự nhiên vang lên khắp lớp học:

``Em là Naomi. Em là họ hàng xa của Suzuki-onee-chan ạ.''

Phù\dots\ May thế.

Ba chúng tôi đồng loạt thở phào nhẹ nhõm, mồ hôi mẹ mồ hôi con tuôn ra như tắm. Thật may là bộ não 1300 tuổi của Nao nhạy bén để ứng phó với tình huống này. Chứ nếu con bé mà khai thật thì chắc ngày mai cả lớp này xin chuyển trường hết mất.

Tiếng chuông reo báo hiệu giờ vào lớp vang lên, cắt ngang sự nhốn nháo. Mọi người miễn cưỡng quay về chỗ ngồi, nhưng những ánh mắt vẫn liên tục liếc về phía cô bé mắt xanh lục đang ôm hộp sữa dâu.

Cánh cửa lớp trượt mở. Yuko-sensei ôm tập giáo án bước vào. Với tư cách là một người lớn từng biết rõ về những sự kiện kỳ lạ, cô luôn giữ được vẻ nghiêm nghị điềm tĩnh.

``Cả lớp trật tự. Bắt đầu tiết sinh hoạt nào--''

Câu nói của cô đứt nghẹn ngay giữa chừng. Ánh mắt cô dừng lại ở chỗ ngồi ngay cạnh cửa sổ, nơi Nao đang bình thản đưa mắt nhìn cô giáo. Tập giáo án trên tay Yuko-sensei rơi cái rạch xuống sàn nhà. Miệng cô há hốc, đôi mắt vốn sắc sảo thường ngày giờ mở to hết cỡ. Cô đưa tay lên tháo mắt kính ra, chùi chùi vào áo rồi đeo lại, nhưng sự thật vẫn không thay đổi: một oán linh ngàn năm mà cô từng nghe nói tới nay đang ngồi chễm chệ trong lớp học của mình, sống sờ sờ bằng xương bằng thịt.

Thế nhưng, trái với phản ứng hoảng loạn của đám học trò, Yuko-sensei chỉ mất đúng ba giây để xóc lại tinh thần. Ngay sau đó, cô hít một hơi thật sâu, từ từ cúi xuống nhặt tập giáo án lên, phủi bụi rồi hắng giọng một cái rõ to. Ánh mắt cô lập tức lấy lại vẻ sắc lạnh thường ngày, phong thái của một nhà giáo mẫu mực ngay lập tức được tái thiết lập.

``E hèm!'' cô chủ nhiệm gõ thước lên bàn giáo viên. ``Cả lớp, về hết chỗ ngồi! Suzuki-san, cuối giờ ở lại gặp tôi tại phòng giáo viên. Còn bây giờ, mở sách giáo khoa ra!''

Nhìn biểu cảm lấy lại phong độ nhanh như chớp của Yuko-sensei, tôi chỉ biết giấu mặt vào hai bàn tay, thầm than trời. Ngày đầu tuần của nhà Hikawa\dots\ xem ra đã bắt đầu theo một cách không thể nhốn nháo hơn rồi.

\ct

Tiết Sinh Học ngay sau tiết sinh hoạt diễn ra trong một bầu không khí kỳ lạ bậc nhất kể từ khi tôi nhập học. Hầu như chẳng ai tập trung hoàn toàn vào bài giảng của thầy giáo cả. Cứ cách vài phút, lại có những ánh mắt từ khắp lớp len lén liếc về phía chỗ ngồi gần cửa sổ. Thế nhưng, trái với sự lo sợ ban đầu của nhóm Shino hay sự tò mò của những người khác, Nao chỉ đơn thuần ngồi đó. Con bé đung đưa chân, chớp đôi mắt lục bảo tò mò quan sát hình minh họa tế bào trên bảng đen, hoàn toàn không tỏa ra bất cứ thứ oán khí hay sát khí nghẹt thở nào như trong đêm nghi lễ. Một sự tồn tại tĩnh lặng, thuần khiết và vô hại.

Khi tiếng chuông báo hiệu hết tiết vang lên, theo lệnh của Yuko-sensei từ đầu giờ, cả Kuon-kun, tôi, Suzu, cùng hai người nữa là Shino và Sakura cùng nhau rời lớp, hướng thẳng tới phòng giáo viên.

Bên trong căn phòng yên tĩnh với mùi cà phê phảng phất, Yuko-sensei ngồi đan hai tay vào nhau trên bàn làm việc, ánh mắt sắc lẹm lướt qua năm người chúng tôi. ``Được rồi. Bắt đầu từ bộ ba Hikawa. Giải thích cho tôi nghe, tại sao `vấn đề tâm linh' cấp độ cao nhất của Aogami bây giờ lại mặc váy hoa, uống sữa dâu và ngồi trong lớp học của tôi?''

Chúng tôi đành phải kể lại toàn bộ sự việc. Từ buổi sáng thức dậy với sự xuất hiện của Nao, cho tới việc Game đã can thiệp để giữ lại phần dữ liệu thuần khiết của cô bé, và cơ chế cộng hưởng với Suzu để tái sinh dưới tư cách một con người bằng xương bằng thịt. Nghe qua không khác nào một câu chuyện được thêu dệt nên từ khoa học viễn tưởng pha trộn bí ẩn.

Sau khi nghe xong, bầu không khí trong phòng giáo viên chìm vào một sự im lặng kéo dài.

Sakura ôm lấy ngực, khẽ thở phào một hơi dài: ``Ra vậy\dots\ Nghĩa là, em ấy không còn bị mắc kẹt trong nỗi oán hận nữa. Lời nguyền thực sự đã kết thúc rồi.''

Shino đứng cạnh cũng gật gù, dù đôi vai vẫn hơi run rẩy do tàn dư của nỗi sợ: ``Bảo sao mình không hề cảm nhận thấy sự lạnh lẽo hay đáng sợ nào từ con bé cả. Ngược lại\dots\ trông em ấy còn có vẻ gì đó rất ngây thơ.''

Yuko-sensei thở dài một hơi, không rõ là do sự bất lực hay vì nhẹ nhõm: ``Thực sự là một chuyện hoang đường ngoài sức tưởng tượng. Nhưng với tư cách là người đã chứng kiến đủ những thứ bất thường ở thị trấn này, tôi không có lý do gì để bác bỏ. Tôi tin các em nói thật.''

``Naomi-chan tuyệt đối không gây hại tới ai đâu ạ,'' Suzu bước lên một bước, ánh mắt kiên định nhìn thẳng vào giáo viên chủ nhiệm. ``Cô bé chỉ là một đứa trẻ đã bị cướp đi quá nhiều thứ mà đáng lẽ em vốn dĩ phải được nhận. Em tin rằng, hiện tại, con bé chỉ muốn được sống như một người bình thường thôi.''

Lời khẳng định của Suzu như một viên đá thả xuống mặt hồ tĩnh lặng. Cả Shino, Sakura và Yuko-sensei đều nhìn nhau, rồi chậm rãi gật đầu đồng tình. Rõ ràng, những gì Nao thể hiện trong suốt hai tiết học vừa qua đã chứng minh điều đó.

``Tôi hiểu rồi,'' Yuko-sensei đứng lên, vỗ nhẹ vai Suzu. ``Nếu cô bé đã thực sự được tái sinh dưới tư cách một con người, thì bổn phận của chúng ta là bảo vệ cuộc sống mới đó. Các em hãy giữ kín bí mật về thân phận thực sự của con bé trước những học sinh khác để tránh rắc rối không đáng có. Tôi sẽ nói chuyện với hiệu trưởng về việc làm thủ tục giấy tờ tạm trú cho cô bé sau. Còn bây giờ, quay về lớp chuẩn bị cho tiết tiếp theo đi.''

``Vâng ạ!'' Cả năm người chúng tôi đồng thanh đáp rồi rảo bước trở về lớp.

Nhưng vấn đề khó nhằn nhất vẫn đang đợi chúng tôi ở lớp 1-C. Những người không biết gì thì dễ lừa bằng cái mác ``họ hàng xa''. Tuy nhiên, với những người đã trực tiếp đối mặt với oán linh ở rừng Aogami như Nanase, Mari, Azuki, Shion, và nhóm đi lần hai như Touka, Saya, Aku, thì vài câu nói dối qua loa không thể qua mặt được họ.

Vừa bước vào cửa lớp, chúng tôi đã bị nhóm này kéo tọt xuống cuối góc lớp.

``Nói thật đi! Cô bé đó có đúng là cái bóng trắng trong rừng không?'' Nanase gắt gao hỏi, Mari đứng cạnh liên tục gật đầu.

``Cả đôi mắt xanh đó nữa! Không thể nào có sự trùng hợp kỳ lạ thế được!'' Touka tiếp lời.

Suzu liếc nhìn Kuon-kun, rồi hít một hơi, quyết định tung ra bài thuyết phục đã chuẩn bị sẵn: ``Mọi người bình tĩnh đã. Đúng là hình dáng bên ngoài có chút tương đồng\dots\ Nhưng đó chỉ là vì Naomi mang trong mình dòng máu của một gia tộc cổ xưa có liên kết ngoại cảm mạnh mẽ với mình. Đêm đó, chính những ý niệm thuần khiết của con bé đã hóa thành hình dáng đó để bảo vệ mọi người khỏi bóng tối.''

Tôi âm thầm giơ ngón cái cho sự bịa đặt nhưng vô cùng thuyết phục của Suzu.

Nghe đến đây, Shion - cậu bạn luôn giữ vẻ điềm tĩnh - khẽ xoa cằm: ``Nói vậy\dots\ cô bé không phải là oán linh, mà là một linh thể bảo hộ đã mượn hình dạng đó sao?''

``Đúng! Chính xác là như vậy!'' Kuon-kun chộp lấy ngay ý tưởng đó, gật đầu lia lịa. ``Và bây giờ con bé đã tới đây để sống cùng gia đình Hikawa. Nên mọi người hoàn toàn yên tâm, không có bất kỳ lời nguyền nào đeo bám cùng chúng ta cả.''

Nhóm Azuki, Nanase và Touka nhìn nhau một hồi, rồi thở phào nhẹ nhõm. Lời giải thích tuy có phần ``tâm linh'' nhưng lại rất hợp lý với thế giới quan mà họ đã trải qua. Còn Suzune, cô nàng não cá vàng nhất nhóm thì vỗ tay cái bốp: ``Vậy tóm lại con bé là họ hàng của cậu và không cắn người chứ gì? Thế thì lo làm gì cho mệt, để mình ra nựng má nó tiếp đây!''

Thế đấy. Nhờ vào sự linh hoạt của Suzu, sự cả tin (và một chút tự biên tự diễn) của các bạn cùng lớp, quả bom nổ chậm lớn nhất trong ngày đầu tuần đã được gỡ bỏ một cách êm thấm.

Nhịp sống học đường lại tiếp diễn, lớp 1-C dần chấp nhận sự hiện diện của ``người nhà'' nhỏ tuổi này. Và tất nhiên, với một cô bé tám tuổi vừa thoát khỏi sự cô lập ngàn năm, mỗi ngày trôi qua đều mở ra vô vàn những câu chuyện nhỏ thú vị đáng để kể lại.

\newpage

\subsubsection*{\phantomsection\label{ep-ER1}}
\addcontentsline{toc}{subsubsection}{{\sen{-Histoire 1-}} {\caviar Bí mật của ngoại cảm}}
\stpt{-Histoire 1-}\stcap{Bí mật của ngoại cảm}
\begin{center}{\caviar Nhân vật: \emp{korone} \emp{okayu}}\end{center}

\begin{center}
  \pov{-}
\end{center}

Giờ nghỉ trưa tại trường trung học Aogami luôn mang đến một nhịp điệu yên ả, trái ngược hoàn toàn với sự nhốn nháo của buổi sáng sớm.

Dưới bóng râm của một gốc cây cổ thụ sau khu tập thể chất,  Yuka đang nằm dài trên chiếc ghế băng bằng gỗ, cuộn người lại như một con mèo tận hưởng ánh nắng mùa hè lọt qua kẽ lá. Gương mặt cô hiện lên vẻ hạnh phúc tột độ của một người có thể dễ dàng chìm vào giấc ngủ ở bất cứ đâu.

Akane ngồi ngay bên cạnh, nhẹ nhàng thu dọn hộp cơm trưa của cả hai. Dù đã quen với cái nết ăn xong là lăn ra ngủ của cô bạn thân, Akane vẫn không giấu được một chút suy tư in hằn trên nét mặt vui vẻ thường ngày. Cô khẽ đưa tay lay nhẹ vai cô bạn.

``Yuka này\dots\ Tôi có chuyện muốn hỏi bà đây.''

Yuka khẽ nhíu mày, mở một nửa đôi mắt ngái ngủ ra nhìn cô bạn thời thơ ấu, miệng ngáp một cái dài: ``Chuyện gì thế Akane? Tôi đang ngủ ngon mà\~{}''

``Chuyện về cô bé lúc sáng ấy,'' Akane ngập ngừng, giọng trầm xuống. ``Bà nói `cơ thể mới' là sao? Đừng nói với tôi con bé thực sự là\dots\ cái bóng mờ mà tôi từng thấy lúc đi khảo sát nhé?''

Nàng bờm Mèo khẽ cựa mình, chống tay ngồi dậy. Sự tĩnh lặng hiếm hoi thay thế cho dáng vẻ ngái ngủ thường ngày. Với năng lực ngoại cảm mạnh mẽ bẩm sinh, cô có thể nhìn thấy và cảm nhận được những thứ mà người bình thường không thể. Đối với cô, việc thấy một linh hồn vất vưởng hay một tàn dư oán khí vốn là chuyện quá đỗi bình thường và chẳng đáng để tâm.

``Thì đúng là con bé đó mà,'' Yuka trả lời tỉnh bơ, đưa tay vươn vai. ``Nhưng bà không cần phải làm cái bộ mặt lo lắng như người mất sổ gạo thế đâu. Nhăn nhó hoài là hết xinh đấy.''

Akane cau mày, nắm lấy ống tay áo của cô bạn: ``Nhưng mà\dots\ sao bà lại tỏ ra bình thản như thế? Dù hình bóng lúc đó rất mờ nhạt, nhưng linh hồn đó từng tỏa ra một bầu không khí rất ngột ngạt cơ mà. Bà không sợ con bé sẽ gây nguy hiểm cho mọi người sao?''

Nhìn nét mặt nghiêm trọng của người bạn luôn hết lòng vì mình, Yuka bật cười khanh khách. Cô vươn tay, nhéo nhẹ vào má Akane.

``Đau, đau, đau! Bà làm gì thế!?''

``Đồ ngốc này,'' Yuka nhe răng cười, giọng nói vang lên nhẹ bẫng như gió. ``Bà chỉ nhìn thấy cái vỏ bọc mờ nhạt lúc con bé còn lang thang, chứ đâu có nhìn thấu được bản chất bên trong giống tôi. Sáng nay, khi chạm vào lọn tóc của Naomi, bà biết tôi cảm nhận được gì không?''

Akane xoa xoa má, tò mò nhìn bạn: ``Bà cảm nhận được gì?''

``Một khoảng không trong suốt. Không oán hận, không nuối tiếc, cũng chẳng còn chút u buồn nào của quá khứ cả,'' cô nàng bờm Mèo ngước mắt nhìn lên những tán lá xanh rì. ``Nó thuần khiết đến mức giống như một tờ giấy trắng vừa được tẩy sạch vậy. Hơn nữa, nếu con bé thực sự nguy hiểm, thì đám người nhà Hikawa đã chẳng dung túng và tự ý đem vào lớp như thế đâu. Bà cứ tin ở họ đi.''

Lời giải thích thong dong của Yuka khiến những khúc mắc trong lòng cô nàng tóc nâu dần được gỡ bỏ. Trước khi quen Yuka, cô vốn là một người ít khi biểu lộ cảm xúc và hiếm khi mỉm cười. Chính nhờ có Yuka - người luôn thấu hiểu và sống thật với chính mình - mà Akane mới tìm thấy niềm vui trong cuộc sống thường nhật. Giờ đây, khi nghe Yuka khẳng định một cách chắc nịch như vậy, cô cũng không còn lý do gì để nghi ngờ nữa.

``Vậy ra, đó thực sự là một phép màu nhỉ?'' nàng bờm Chó lẩm bẩm, khóe môi khẽ nhếch lên thành một nụ cười an tâm.

``Đúng thế\~{} Phép màu của thời đại này đấy,'' Yuka ngáp thêm một cái nữa, rồi lại tự nhiên ngả đầu xuống đùi Akane. ``Nên là\dots\ đừng lo bò trắng răng nữa. Giờ thì cho tôi mượn đùi ngủ nốt mười lăm phút đi\dots\ Zzz\dots''

Chưa đầy năm giây sau, nhịp thở của cô nàng bờm Mèo đã trở nên đều đặn. Akane cúi xuống nhìn gương mặt ngủ say sưa của bạn mình, bất giác bật cười thành tiếng. Cô tìm thấy niềm hạnh phúc đơn giản chỉ bằng việc đánh thức và chăm sóc cho giấc ngủ của cô bạn thân này. Khẽ vuốt lại mái tóc rối của Yuka, trong lòng Akane tràn ngập cảm giác bình yên. Phải rồi, miễn là vẫn còn những khoảnh khắc bình dị như thế này, thì dù là oán linh tái sinh hay phép màu nào đi nữa, thế giới của họ vẫn sẽ ổn thôi.

Cùng lúc đó, từ bậu cửa sổ hành lang tầng hai nhìn thẳng xuống khu vực gốc cây cổ thụ, Naomi đang nhón chân bám tay vào lan can. Đôi mắt lục bảo của em tĩnh lặng quan sát hai nữ sinh lớp 1-C.

Suzuki bước tới từ phía sau, áp một hộp sữa dâu nhỏ mát lạnh vào má em khiến con bé khẽ giật mình rụt cổ lại.

``\vie{Em đang nhìn gì thế, Naomi-chan?}'' cô nàng mắt hai màu tò mò hỏi.

Cô bé tóc tết nhận lấy hộp sữa, ngước lên mỉm cười: ``\vie{Em đang nhìn một người có thể nhìn thấu em. Chị ấy\dots\ có một tâm hồn rất đẹp và tự do, giống như đám mây mùa hè vậy.}''

Suzu nhìn theo hướng mắt của con bé, bật cười khi thấy Yuka đang ngủ ngon lành trên đùi Akane dưới gốc cây. ``\vie{Phải rồi, Yuka-san luôn có một cách nhìn thế giới rất đặc biệt mà.}''

Naomi cắm ống hút vào hộp sữa, nhấp một ngụm nhỏ, đôi mắt lấp lánh phản chiếu ánh mặt trời rực rỡ. Thế giới này không chỉ có bảng quảng cáo LED lấp lánh hay xe scooter phát nhạc ồn ào, mà còn có những tâm hồn ấm áp sẵn sàng dang tay đón nhận em. Và điều đó, với cô bé tám tuổi vừa tỉnh giấc từ ngàn năm bóng tối, mới là phép màu tuyệt diệu nhất.

\newpage

\subsubsection*{\phantomsection\label{ep-ER2}}
\addcontentsline{toc}{subsubsection}{{\sen{-Histoire 2-}} {\caviar Lời chào từ phòng phát thanh}}
\stpt{-Histoire 2-}\stcap{Lời chào từ phòng phát thanh}
\begin{center}{\caviar Nhân vật: \emp{towa} \emp{kanata}}\end{center}

Bên trong phòng phát thanh của trường trung học Aogami, Kana đang bận rộn sắp xếp lại những đĩa CD nhạc nền cho buổi phát sóng buổi trưa. Mái tóc hai bím xoáy của cô khẽ đung đưa theo từng nhịp điệu của bài hát đang lẩm nhẩm trong miệng. Trái ngược với vẻ năng nổ, hoạt bát của cô nàng, cô bạn thời thơ ấu Kanade lại ngồi ngay ngắn trên chiếc ghế cạnh micro, cẩn thận lật xem từng trang kịch bản với vẻ mặt đầy nghiêm túc.

``Này Kanade-chan, bà nghĩ hôm nay chúng ta nên thêm một góc phỏng vấn đặc biệt không?'' Kana đột ngột quay sang, đôi mắt sáng rực lên sự tinh nghịch. ``Chủ đề là `Nhân vật bí ẩn mới xuất hiện ở lớp 1-C'. Tôi nghe nói nhà Hikawa vừa mang một cô bé mới tám tuổi đến trường đấy. Nếu tớ vác micro xuống phỏng vấn cô bé đó, rating đài phát thanh trường mình chắc chắn sẽ nổ tung!''

Cô bạn Thiên sứ giật mình, vội vàng xua tay can ngăn: ``Không được đâu Kana-chan! Làm thế sẽ khiến em ấy hoảng sợ mất.''

``Sợ gì chứ? Tôi là một người dẫn chương trình vô cùng thân thiện mà!'' Kana bĩu môi, hai tay chống nạnh.

Kanade khẽ thở dài, nụ cười dịu dàng xen lẫn chút lo lắng hiện lên trên gương mặt. Dù có tính cách kín đáo, nhưng trái tim nhân hậu và sự nhạy cảm của cô lại giúp cô nhìn nhận sự việc sâu sắc hơn bề ngoài.

``Lúc nãy khi đi ngang qua hành lang, tôi đã vô tình chạm mắt với cô bé tên Naomi-chan ấy,'' cô trầm giọng, nhớ lại ánh mắt lục bảo trong veo nhưng lại ẩn chứa một độ sâu tĩnh lặng kỳ lạ của cô bé tám tuổi. ``Đó không phải là ánh mắt của một đứa trẻ bình thường đâu, Kana-chan. Tôi có cảm giác\dots\ em ấy như vừa trải qua một giấc ngủ rất dài và chỉ mới bắt đầu học cách nhìn ngắm thế giới này vậy. Nếu chúng ta đột ngột kéo em ấy vào sự chú ý của đám đông, em ấy có thể sẽ thu mình lại đó.''

Nghe người bạn thân phân tích, Kana chớp chớp mắt, rồi thở hắt ra một hơi đầu hàng. Bề ngoài luôn tỏ ra mạnh mẽ và thích làm trung tâm, nhưng bên trong cô lại là một người vô cùng nhạy cảm và dễ thấu hiểu người khác.

``Được rồi, được rồi, thua bà luôn!'' cô nàng Ác ma đi tới, dùng cuốn kịch bản gõ nhẹ lên đầu Kanade một cái. ``Bà lúc nào cũng nhân hậu như một thiên thần ấy. Thế thì thôi, dẹp ý tưởng phỏng vấn vậy. Nhưng bù lại, tôi sẽ mở một bài hát thật vui nhộn, coi như là lời chào mừng không chính thức tới thành viên nhỏ tuổi nhất của trường mình nhé!''

Kanade mỉm cười rạng rỡ, gật đầu: ``Ừm, ý kiến tuyệt lắm!''

Vài phút sau, chiếc loa phóng thanh gắn trên góc sân trường bắt đầu vang lên giai điệu trong trẻo của một bản nhạc pop mùa hè. Giọng nói lưu loát, dõng dạc của Kana cất lên, theo sau đó là giọng nói mềm mại, êm ái của cô nàng Thiên sứ - người nhờ những lần làm khách mời mà đã tự tin hơn hẳn khi nói trước đám đông.

Ở dưới sân trường, trong lúc đang đứng chờ Suzuki đi mua đồ uống từ máy bán hàng tự động, Kaigou đang khoanh tay tựa lưng vào tường. Đứng ngay bên cạnh cô là Naomi.

Cô bé tóc tết ngước nhìn chiếc hộp sắt vuông vức treo trên cao đang phát ra âm thanh, đôi mắt lục bảo mở to đầy kinh ngạc.

``Kaigou-onee-chan,'' cô bé giật giật vạt áo của Kaigou, ngón tay nhỏ nhắn chỉ lên chiếc loa. ``\vie{Tại sao trong cái hộp sắt đó lại có giọng nói của con người? Chẳng lẽ\dots\ lại có linh hồn nào bị giam trong đó sao?}''

Cô nàng tóc đuôi ngựa phì cười, xoa đầu cô bé: ``\vie{Không có linh hồn nào bị giam cầm đâu, Nao ạ. Đó là hệ thống loa phát thanh của trường. Có hai chị nữ sinh đang ngồi ở một căn phòng khác nói vào một thiết bị tên là micro, và âm thanh sẽ được truyền qua dây điện để phát ra ngoài này.}''

Naomi chớp mắt, dường như vẫn chưa thể tiêu hóa hết khái niệm về ``dây điện'' và ``micro'', nhưng em nhanh chóng thả lỏng đôi vai nhỏ bé. Em khép hờ mắt lại, lặng lẽ lắng nghe sự hòa quyện giữa giọng nói năng động của Kana và chất giọng êm đềm của Kanade.

``\vie{Không có sự oán hận\dots\ cũng không có tiếng khóc than,}'' cô bé thì thầm, khóe môi khẽ cong lên. ``\vie{Âm thanh phát ra từ đó\dots\ ấm áp giống như một lời chúc phúc vậy.}''

Kaigou nhìn cô bé, khẽ mỉm cười. ``Ừ. Đó chính là giai điệu của thế giới người sống đấy.''

\newpage

\subsubsection*{\phantomsection\label{ep-ER3}}
\addcontentsline{toc}{subsubsection}{{\sen{-Histoire 3-}} {\caviar Sự ấm áp từ những vật vô tri}}
\stpt{-Histoire 3-}\stcap{Sự ấm áp từ những vật vô tri}
\begin{center}{\caviar Nhân vật: \emp{flare} \emp{luna} \emp{roboco}}\end{center}

Khi tiếng chuông báo hiệu kết thúc giờ học buổi chiều vang lên, lớp 1-C nhanh chóng chìm vào không khí rộn rã thường ngày. Kuon đang lúi húi nhét sách vở vào cặp, còn Naomi thì ngoan ngoãn ngồi chờ trên ghế, đôi mắt tò mò nhìn mọi người qua lại.

Đúng lúc đó, Hotaru lao tới bàn của hai người với tốc độ của một cơn lốclốc, hớn hở sà tới cô bé mắt xanh đang đung đưa đó.

``Naomi-chan! Chị có cái này cho em này!'' Hotaru reo lên, chìa ra một con gấu bông nhỏ xíu có gắn móc khóa, lông mềm mượt và cái nơ hồng xinh xắn. ``Hàng giới hạn đấy nhé! Tặng cho thành viên dễ thương nhất lớp mình!''

Ngay khi Hotaru định lao tới ôm chầm lấy Nao, một bàn tay thon thả nhưng đầy lực cản đã túm lấy cổ áo cô nàng kéo lại. Mitsuki bước tới, mỉm cười dịu dàng, nhưng sát khí vô hình lại khiến Hotaru phải đổ mồ hôi lạnh.

``Hotaru-chan, bớt vồ vập lại đi. Cậu đang làm em ấy sợ đấy na, '' chất giọng của cô nàng tóc hồng êm ái như kẹo bông, nhưng bên trong lại đầy uy lực của một người vốn quen làm chủ tình hình.

``Gì chứ\~{}? Tôi chỉ muốn tặng quà thôi mà\dots'' Hotaru phồng má cãi lại, nhưng vẫn ngoan ngoãn lùi lại một bước.

Từ phía sau, Yuki khoanh tay bước tới, nhếch mép mỉa mai: ``Khỏi giả bộ đi Hotaru. Ai trong lớp này mà chẳng biết phòng cậu chất đầy thú bông tới mức không còn chỗ ngủ? Cậu chỉ tìm cớ để lôi đồ chơi ra khoe thôi đúng không?''

``N-Nói bậy! Tôi chỉ thấy hợp với Naomi-chan thôi!'' Hotaru đỏ bừng mặt, cuống quýt xua tay. Bộ dạng lúng túng khi bị khen ngợi hay bóc mẽ của cô khiến cả Mitsuki và Yuki bật cười khúc khích.

Ở giữa màn đấu khẩu trêu chọc đó, Naomi từ từ đưa cả hai tay ra đón lấy con gấu bông nhỏ. Đôi mắt lục bảo mở to, chăm chú quan sát món đồ vật trong tay. Nó hoàn toàn vô tri vô giác, không được yểm bất kỳ bùa chú nào, cũng không chứa đựng pháp lực bảo vệ. Chỉ là một món đồ chơi bình thường có thể bắt gặp ở bất kỳ cửa hiệu nào.

``\vie{Cái này\dots\ không dùng để trấn yểm sao ạ?}'' Naomi ngước lên nhìn Kuon, giọng nói mang theo sự thắc mắc thuần túy của một người đến từ thời đại mà mọi hình nhân đều mang mục đích tâm linh.

Kuon bật cười, đưa tay xoa nhẹ mái tóc tết của cô bé: ``\vie{Không đâu. Ở thời đại này, người ta tạo ra chúng không phải để xua đuổi tà ma, mà là để mang lại niềm vui. Người ta tặng nó cho nhau để thể hiện sự yêu mến đấy.}''

Naomi chớp mắt, rồi từ từ áp con gấu bông nhỏ vào má. Lớp vải bông mềm mại chạm vào da, truyền tải một cảm giác ấm áp hoàn toàn khác biệt với những pháp khí lạnh lẽo mà em từng biết trong quá khứ. Nó chứa đựng một ý niệm đơn giản mà thuần khiết của Hotaru.

``\vie{Vậy ra, ở thời đại này, người ta dùng sự yêu thương để bảo vệ lẫn nhau,}'' cô bé thì thầm, ôm chặt lấy món quà nhỏ vào ngực, khóe môi vẽ lên một nụ cười rạng rỡ. ``\vie{Em cảm ơn chị, Hotaru-onee-chan.}''

Nhìn nụ cười chân thành của con bé, mặt Hotaru lại càng đỏ hơn. Cô nàng đưa tay gãi đầu, cười trừ ngượng ngùng, trong khi Mitsuki và Yuki trao nhau một ánh nhìn ấm áp. Giữa thế giới hiện đại đầy rẫy những thiết bị vô tri, những món đồ nhỏ bé này vẫn có thể chứa đựng một phép màu\dots\ phép màu của tình người.

\newpage

\subsubsection*{\phantomsection\label{ep-ER4}}
\addcontentsline{toc}{subsubsection}{{\sen{-Histoire 4-}} {\caviar Phù thủy vĩ đại và Oán linh ngàn năm}}
\stpt{-Histoire 4-}\stcap{Phù thủy vĩ đại và Oán linh ngàn năm}
\begin{center}{\caviar Nhân vật: \emp{shion} \emp{aqua}}\end{center}

Ngày hôm sau, trong giờ ra chơi, góc cuối lớp 1-C phát ra những tràng cười giòn giã. Hay nói đúng hơn, chỉ có một mình Aku đang cố gắng nhịn cười trước mặt Shion.

``Đừng có cười! Đây là Ma trận Bóng tối mà ta vừa thức tỉnh đêm qua đấy!'' cô nàng Phù thủy tự xưng bực dọc đập tay xuống bàn. Cô chỉ tay vào trang vở chi chít những vòng tròn ma thuật được vẽ bằng bút bi đỏ đen đan xen. ``Chỉ cần truyền một ít ma năng vào đây, ta có thể triệu hồi một tinh linh hắc ám từ thế giới ngầm!''

Aku một tay ôm bụng, tay kia lau giọt nước mắt rỉ ra nơi khóe mắt: ``Thôi đi Shion-chan\dots\ Bà lại thức đêm cày anime dị giới nữa chứ gì? Trông cái vòng tròn kia y hệt cái logo của hãng nước tăng lực bà hay uống luôn ấy.''

``T-Ta không có chép! Đây là ấn chú chính hiệu!'' Shion đỏ mặt tía tai cãi lại.

Đúng lúc đó, Naomi đang nắm tay Suzuki đi ngang qua hành lang lớp. Đôi mắt lục bảo của cô bé lướt qua trang vở đang mở tung trên bàn của cô nàng tóc xám-trắng. Đột nhiên, cô bé buông tay Suzuki ra, chậm rãi bước tới gần bàn của hai người.

Thấy ``người nhà Hikawa'' tiến lại gần, cả Shion và Aku lập tức cứng đơ người. Cô hầu gái rụt rè vội vàng nuốt nước bọt, ngồi nép sát vào ghế, còn cô Phù thủy thì run rẩy định giấu quyển vở đi. Thế nhưng, Naomi lại rướn người lên, chăm chú nhìn vào những vòng tròn ma thuật mà Shion tự vẽ.

Cô bé tám tuổi nhíu mày, ngón tay nhỏ xíu chỉ vào một góc của trận pháp.

``Chỗ này\dots\ vẽ ngược rồi ạ,'' Naomi cất giọng đều đều, nghiêm túc chỉ ra lỗi sai. ``Nếu chị muốn viết chữ Phạn để kết nối âm giới, nét này phải móc sang trái. Còn ở trung tâm\dots\ nếu ghép các ký tự này lại, thay vì gọi tinh linh, nó sẽ biến thành trận pháp thu hút ngạ quỷ.''

Cả Shion và Aku há hốc mồm. Bầu không khí tĩnh lặng đến mức có thể nghe tiếng gió thổi ngoài cửa sổ.

``Hơn nữa, bút bi đỏ không thể làm vật dẫn ma lực được. Chị phải dùng chu sa trộn với một chút máu gà trống để vẽ lên giấy bồi,'' Naomi nghiêng đầu, đôi mắt trong veo như chứa đựng cả một kho tàng kiến thức tâm linh cổ xưa. Cô bé chìa tay ra: ``Chị có muốn em dạy cách vẽ trận pháp trừ tà của thời đại em không? Hiệu nghiệm lắm đó.''

Bị một chuyên gia ``hàng xịn'' với kinh nghiệm hơn một thiên niên kỷ bắt bẻ ngay trên mặt trận tâm linh, niềm kiêu hãnh của ``Phù thủy vĩ đại'' vỡ vụn ngay lập tức. Mặt nàng Phù thủy đỏ bừng từ cổ lên tận mang tai. Cô nàng luống cuống gấp sầm quyển vở lại, giấu ngay xuống ngăn bàn.

``K-Không cần đâu! Phù thủy vĩ đại như ta tự có cách riêng của mình! Khác hệ, khác hệ ma thuật, hiểu không!?'' Shion vội vàng xua tay, lắp bắp lùi lại phía sau.

Aku ngồi bên cạnh, sau một vài giây sững sờ, liền không nhịn được mà bật cười phá lên. Tiếng cười của cô nàng vang vọng khắp góc lớp, mặc cho Shion bên cạnh đang xấu hổ đến mức muốn đào một cái hố để chui xuống.

Naomi chớp mắt, khó hiểu nhìn phản ứng của hai nữ sinh. Cô bé ngoái đầu lại nhìn Suzuki, người cũng đang bụm miệng cười nắc nẻ.

``\textit{Thế giới hiện đại thật kỳ lạ\dots}'' Naomi thầm nghĩ, ngoan ngoãn nắm lấy tay Suzu để bước tiếp. Dù không hiểu hết khái niệm ``khác hệ ma thuật'' mà Shion vừa nói, nhưng tiếng cười sảng khoái của Aku khiến em cảm thấy lớp học này thực sự rất thú vị. Sẽ chẳng có con ngạ quỷ nào bị triệu hồi cả, vì nơi đây vốn dĩ đã ngập tràn ánh sáng và tiếng cười rồi.

\newpage

\subsubsection*{\phantomsection\label{ep-ER5}}
\addcontentsline{toc}{subsubsection}{{\sen{-Histoire 5-}} {\caviar Nỗi sợ hãi mang màu lá phong}}
\stpt{-Histoire 5-}\stcap{Nỗi sợ hãi mang màu lá phong}
\begin{center}{\caviar Nhân vật: \emp{ayame} \emp{watame}}\end{center}

Trên hành lang dẫn ra sân thể dục, Ayame liên tục ngoái đầu nhìn trước ngó sau, dáng vẻ lấm lét khác hẳn với sự tự tin thường ngày. Vốn dĩ, cô nàng tóc đỏ luôn nổi tiếng là người mạnh dạn, hay trêu chọc và ít khi sợ hãi thứ gì. Nhưng đó là trước khi cô có mặt tại đêm nghi lễ trong khu rừng Aogami.

``Cậu sao vậy Ayame? Cứ đi một đoạn lại giật mình thế?'' Wata nghiêng đầu hỏi, mái tóc xoăn mềm mại của cô nàng khẽ đung đưa.

``T-Tôi không sao! Chỉ là\dots\ tôi đang tập bài thể dục giãn cơ cổ thôi!'' Ayame vội vàng đánh trống lảng, ngẩng mặt lên trời vờ như đang khởi động.

Wata chớp mắt, dường như không mấy tin tưởng lời bào chữa đó. Đúng lúc này, từ phía đầu hành lang bên kia, Kuon và Naomi đang thong thả đi tới. Vừa nhìn thấy cô bé tóc tết với đôi mắt lục bảo quen thuộc, Ayame lập tức đông cứng người, lùi lại nấp ra sau lưng cô bạn.

``K-Không được rồi, Wata-chan\dots\ Cậu ra đỡ đòn cho tôi đi. Em ấy tới kìa\dots'' cô nàng tóc trắng run rẩy bám lấy vai cô bạn, giọng nói mếu máo như sắp khóc. Dù đã được giải thích rằng Naomi là họ hàng của Suzu và không còn là ác linh, nhưng tàn dư ký ức từ đêm đó vẫn ám ảnh cô nàng.

Trái ngược với vẻ hoảng loạn của cô bạn thân, Wata lại sáng rực hai mắt. Tính cách ngây thơ và bản năng yêu thích những thứ nhỏ nhắn, mềm mại lập tức trỗi dậy. Cô nàng nhẹ nhàng gỡ tay Ayame ra, lon ton chạy tới trước mặt Naomi.

``Chào em, Naomi-chan!'' cô nàng tóc vàng hớn hở cúi xuống, đôi mắt lấp lánh nụ cười. Cô nàng khẽ vuốt ve bím tóc của Naomi. ``Oaa\~{} Tóc em mềm thật đấy, cứ như len cừu non vậy! Chị ôm em một cái được không?''

Naomi chớp mắt, rồi ngoan ngoãn dang hai tay ra. Wata lập tức ôm chầm lấy cô bé, cọ cọ má vào đôi má phúng phính của cô với vẻ vô cùng hạnh phúc. Ở đằng xa, Ayame đang há hốc mồm, mồ hôi lạnh chảy ròng ròng trên trán. Trong mắt cô, cô nàng Hội phó đang tự mình bước vào ``vùng đất tử thần''.

``Ayame-chan, cậu cũng lại đây ôm thử xem!'' Wata vẫy tay gọi. ``Người em ấy ấm áp và thơm mùi sữa dâu lắm này! Mình đảm bảo không có gì đáng sợ đâu!''

Bị gọi đích danh, Hội trưởng nuốt nước bọt cái ực, rụt rè tiến lên từng bước một như người đang đi trên bãi mìn. Cô nàng nhắm tịt mắt, chìa một ngón tay ra, chọc chọc nhẹ vào vai Naomi như thể đang kiểm tra xem cô bé có tan biến thành sương mù hắc ám không.

Cảm nhận được cái chạm rụt rè đó, Naomi ngoái đầu lại. Đôi mắt lục bảo tĩnh lặng nhìn sâu vào đôi mắt đang mở to đầy phòng thủ của Ayame. Khi còn là một thực thể tâm linh, em từng đọc được nỗi sợ hãi trong tâm trí cô Hội trưởng này. Và bây giờ, sự sợ hãi ấy vẫn còn, nhưng nó không đi kèm với oán ghét hay thù hằn, mà chỉ là một sự rụt rè vô cùng đáng yêu.

Khẽ nhếch khóe môi, Naomi từ từ đưa bàn tay nhỏ bé của mình ra, chạm nhẹ vào những lọn tóc đỏ rực của Ayame.

``Á!'' Ayame giật thót mình, nhắm nghiền mắt lại chờ đợi một ``lời nguyền'' giáng xuống.

Thế nhưng, thứ vang lên lại là một chất giọng trong trẻo, mang theo một chút hồn nhiên hiếm hoi của đứa trẻ tám tuổi: ``Tóc của chị có màu rất đẹp. Giống hệt như màu của lá phong mùa thu rụng ngoài hiên đền hồi trước vậy. Ấm áp lắm ạ\~{}''

Hội trưởng chớp mắt, sững sờ mở mắt ra. Trông thấy nụ cười thuần khiết không chút tạp niệm của Naomi, nỗi sợ hãi vô hình đè nặng trong lòng cô bỗng chốc tan biến như sương mai dưới nắng. Căng thẳng tiêu tan, cô nàng đột nhiên che miệng, bật ra tiếng cười khúc khích đặc trưng. Nụ cười rạng rỡ của cô nàng tỏa sáng cả dãy hành lang.

``Nghe thấy chưa Wata-chan? Em ấy vừa khen tôi đấy! Tóc màu lá phong cơ đấy!'' Ayame tự hào vuốt tóc, hoàn toàn quên mất việc mình vừa nấp sau lưng bạn thân sợ hãi thế nào.

Wata lắc đầu mỉm cười: ``Mình đã bảo rồi mà. Hội trưởng cứ hay lo xa quá.''

Kuon đứng cạnh khoanh tay thở dài, không khỏi buồn cười trước sự chuyển biến thái độ nhanh như chớp của nàng tóc đỏ. Còn Naomi thì tựa đầu vào vai Wata, lẳng lặng khắc ghi tiếng cười lanh lảnh của Ayame vào trong trí nhớ. Con người ở thời đại này\dots\ thật mong manh, nhát gan, nhưng đôi khi, chính sự ngốc nghếch chân thành của họ lại là thứ bùa chú sưởi ấm tâm hồn hữu hiệu nhất.

\newpage

\subsubsection*{\phantomsection\label{ep-ER6}}
\addcontentsline{toc}{subsubsection}{{\sen{-Histoire 6-}} {\caviar Nữ chiến binh và Máy phát thanh}}
\stpt{-Histoire 6-}\stcap{Nữ chiến binh và Máy phát thanh}
\begin{center}{\caviar Nhân vật: \emp{polka} \emp{botan}}\end{center}

Honoka là một quả bom năng lượng không bao giờ cạn của lớp 1-C. Và dĩ nhiên, một người luôn thích nhúng tay vào mọi rắc rối để khuấy động bầu không khí như cô sẽ không đời nào bỏ qua cơ hội được vây quanh ``thành viên nhí'' mới của lớp.

``Naomi-chan, Naomi-chan! Nhìn này, chị có thể xoay ba cuốn sách trên một ngón tay đấy!'' cô nàng bờm Cáo sa mạc ríu rít không ngừng, múa may quay cuồng trước mặt Naomi. ``Hoặc là\dots\ em muốn nghe chuyện về thầy dạy Toán lớp kế bên bị vấp vỏ chuối ngã đập mặt sáng nay không? Đảm bảo buồn cười nổ bụng luôn!''

Cô bé tám tuổi ngồi im trên ghế, đôi mắt lục bảo chớp liên tục không ngừng. Không phải vì em đang hào hứng, mà vì tốc độ bắn chữ và chuyển động của Honoka quá nhanh khiến một bộ não vốn đã ngủ yên ngàn năm chưa kịp xử lý lượng thông tin khổng lồ ấy.

Từ dãy bàn phía cuối, Shiina đang gác chân lên bàn, mắt dán vào màn hình điện thoại chơi game. Vị ``đại tỷ'' được mệnh danh là học sinh cá biệt mạnh nhất Aogami này khẽ chậc lưỡi.

``Tôi nói này, đừng có làm rộn lên thế Honoka, cậu làm con bé chóng mặt rồi kìa,'' cô lười biếng cất giọng, âm vực trầm ổn mang theo một chút uy lực khiến những người xung quanh phải nể sợ.

``Tớ chỉ muốn em ấy hòa nhập nhanh thôi mà!'' Honoka phồng má quay lại cãi. ``Cậu nhìn xem, nãy giờ Naomi-chan chớp mắt liên tục, chắc chắn là đang rất hứng thú với mấy trò của tớ!''

Cô nàng bờm Sư tử thở hắt ra, tắt màn hình điện thoại rồi nhét vào túi áo, rồi chậm rãi bước tới. Không nói không rằng, Shiina đưa một tay ra, nhẹ nhàng túm lấy cổ áo phía sau của cô bạn bát nháo và nhấc bổng cô nàng sang một bên như xách một con mèo con.

``Nhẹ thôi nhẹ thôi! Cổ áo tớ rách bây giờ!'' Honoka chới với giãy nảy, nhưng sức mạnh của nữ chiến binh đường phố này là không thể bàn cãi.

Sau khi ``dọn dẹp'' xong chướng ngại vật ồn ào, cô nàng bờm Sư tử kéo ghế ngồi phịch xuống đối diện Naomi. Khác với vẻ ngoài có phần bặm trợn và sát khí ngầm của một người hay đánh lộn, Shiina lúc này lại tỏa ra một sự điềm tĩnh kỳ lạ. Cô thò tay vào túi áo, lấy ra một chiếc kẹo mút vị dâu tây, bóc vỏ rồi đưa ra trước mặt cô bé.

``Đừng để ý cái máy phát thanh chạy bằng cơm kia. Lắm mồm nhưng không có ác ý đâu,'' Shiina khẽ nhếch mép cười, một nụ cười cực kỳ ngầu. ``Ăn kẹo không?''

Naomi rụt rè vươn tay nhận lấy chiếc kẹo. Em chậm rãi đưa vào miệng. Vị ngọt lịm của dâu tây nhân tạo ngay lập tức lan tỏa trên đầu lưỡi, khiến đôi mắt lục bảo vốn đang choáng váng bỗng chốc sáng rực lên.

``Ngọt quá ạ\~{}'' Naomi thì thầm, hai má ửng hồng.

``Thấy chưa? Trẻ con thì cứ phải dùng đồ ngọt mới tĩnh tâm lại được,'' Shiina khoanh tay, đắc ý gật gù, trong khi Honoka ở phía sau đang giậm chân bình bịch.

``Aaa! Shiina ăn gian! Tớ cũng định cho em ấy xem bộ sưu tập kẹo bông gòn của tớ cơ mà!'' nàng bờm Cáo sa mạc gào lên phản đối.

Mặc kệ tiếng ồn ào vẫn tiếp tục văng vẳng bên tai, Naomi khẽ ngậm viên kẹo mút. Dưới góc nhìn của cô bé, cô nàng tóc bạch kim trước mặt giống hệt như một chiến binh thực thụ: mang trong mình sát khí mạnh mẽ, nhưng lại biết cách kiểm soát nó hoàn toàn để bảo vệ sự yên bình. Còn cô gái ồn ào kia tuy hơi phiền phức, nhưng lại mang đến nguồn năng lượng vui vẻ có thể xua tan mọi bóng tối. Hai con người trái ngược\dots\ nhưng lại là một sự kết hợp hoàn hảo đến kỳ lạ của thế giới hiện đại.

\newpage

\subsubsection*{\phantomsection\label{ep-ER7}}
\addcontentsline{toc}{subsubsection}{{\sen{-Histoire 7-}} {\caviar Những vì sao và mẩu sô-cô-la}}
\stpt{-Histoire 7-}\stcap{Những vì sao, hộp sữa và mẩu sô-cô-la}
\begin{center}{\caviar Nhân vật: \emp{mio} \emp{subaru} \emp{mel}}\end{center}

Ở một góc khác của lớp 1-C, Miku đang ngồi vắt chân, tay lật giở một cuốn tạp chí thời trang mới nhất. Cô nàng với mái tóc đen dài và cặp (bờm) tai sói ẩn hiện luôn cố gắng xây dựng hình ảnh một người phụ nữ trưởng thành, quý phái.

``Mùa thu năm nay, tông màu trầm sẽ lên ngôi đấy,'' Miku nhận xét với vẻ mặt đầy chiêm nghiệm, trong khi tay kia lén lút bóc một mẩu sô-cô-la nhỏ bỏ vào miệng. ``Ăn sô-cô-la đen không chỉ giúp thư giãn mà còn làm đẹp da nữa. Đó là bí quyết của những quý cô.''

``Miku-chan, cậu lại đọc mấy cái tạp chí lá cải đó nữa à?''  Kaoru vừa nói vừa hì hục tu hết một hộp sữa lớn. Cô nàng thở hắt ra một hơi sảng khoái: ``Muốn đẹp thì phải vận động! Và uống sữa nữa! Tôi nhất định phải cao thêm ít nhất hai xăng-ti-mét nữa trong tháng này!''

Saya, người nãy giờ vẫn im lặng nhìn ra cửa sổ với ánh mắt mơ màng, bỗng cất giọng với nhịp điệu chậm rãi đặc trưng: ``Mọi người biết không\~{} Trên bầu trời lúc này, dù chúng ta không thấy, nhưng chòm sao Thiên Long đang nằm ở vị trí rất đẹp đó\~{} Nó như một con rồng đang bảo vệ cả thị trấn Aogami vậy\~{}''

Đúng lúc đó, Naomi được Suzuki dắt tay đi ngang qua bàn của họ. Cô bé dừng lại, tò mò nhìn nàng bờm Sói đang vội vàng giấu mẩu sô-cô-la đi vì xấu hổ khi bị trẻ con bắt gặp cảnh ``làm điệu''.

``Chị đang ăn gì thế ạ?'' Naomi nghiêng đầu hỏi, đôi mắt lục bảo trong veo.

Miku đỏ mặt, ấp úng: ``À\dots\ Đây là ma lực đen giúp làm đẹp-- Ý chị là, sô-cô-la thôi. Em muốn thử không?''

Cô bẻ một miếng nhỏ đưa cho Naomi. Cô bé nhận lấy, nếm thử rồi khẽ nhăn mặt vì vị đắng, nhưng sau đó lại mỉm cười khi vị ngọt hậu lan tỏa.

Trong lúc đó, cô nàng Thiên văn cúi xuống, dịu dàng xoa đầu Naomi và bắt đầu giảng giải về sự tích các ngôi sao bằng chất giọng như tiếng chuông gió: ``Trên bầu trời đêm ngàn năm trước, chắc hẳn em đã thấy những vì sao rực rỡ hơn bây giờ nhiều nhỉ\~{} Chúng là những linh hồn đang dõi theo chúng ta đó\~{}''

Naomi lắng nghe một cách say sưa. Đối với một người từng sống ở thời đại mà bầu trời đêm là thứ duy nhất soi sáng bóng tối, kiến thức của Saya khiến em cảm thấy vô cùng gần gũi.

``\textit{Chị ấy nói chuyện nghe rất êm tai,}'' Naomi thầm nghĩ, trong khi Kaoru đang hào hứng tặng em thêm một hộp sữa nhỏ. ``\textit{Một người cố gắng làm người lớn, một người năng động như mặt trời, và một người tĩnh lặng như ánh trăng. Thế giới này thực sự có rất nhiều màu sắc thú vị.}''

Cô bé tám tuổi ôm hộp sữa vào lòng, cảm thấy mỗi người bạn mới trong lớp 1-C đều mang đến một dư vị khác nhau, giống như miếng sô-cô-la vừa rồi: ban đầu có chút lạ lẫm, nhưng càng về sau lại càng ngọt ngào.

\newpage

\subsubsection*{\phantomsection\label{ep-ER8}}
\addcontentsline{toc}{subsubsection}{{\sen{-Histoire 8-}} {\caviar Ngọn lửa và Mặt hồ tĩnh lặng}}
\stpt{-Histoire 8-}\stcap{Ngọn lửa và Mặt hồ tĩnh lặng}
\begin{center}{\caviar Nhân vật: \emp{matsuri} \emp{fubuki}}\end{center}

Ở giữa lớp học, nơi không khí luôn sôi động nhất, Hanabi đang đứng chống nạnh, cười rạng rỡ với cô bạn thân ngồi đối diện. Cô nàng đang cố gắng truyền ngọn lửa nhiệt huyết của mình sang cho lớp trưởng Inari.

``Đi mà Inari-chan! Bà chỉ cần đứng lên bục và thông báo về kế hoạch dã ngoại sắp tới thôi!'' Hanabi vỗ mạnh vào vai bạn, giọng nói đầy tính khích lệ. ``Có tôi đứng ngay phía sau rồi, không ai dám làm khó đâu!''

Inari nắm chặt cuốn sổ ghi chép, đôi tai (bờm) cáo của cô khẽ cụp xuống. Dù đã nỗ lực làm lớp trưởng để thay đổi tính cách nhút nhát, nhưng mỗi khi phải đứng trước đám đông, cô vẫn không tránh khỏi cảm giác lo âu. Cô có thói quen đọc tâm trạng người khác quá kỹ, đôi khi điều đó khiến cô trở nên quá nhạy cảm với những ánh mắt xung quanh.

``Nhưng mà Hanabi-chan\dots\ tôi sợ mình sẽ nói vấp, hoặc mọi người sẽ thấy chán mất\dots'' Inari ấp úng, ánh mắt dán chặt vào mặt bàn.

Đúng lúc đó, Naomi đang đi dạo quanh lớp cùng với Kaigou. Cô bé dừng chân lại cạnh bàn của hai người, đôi mắt lục bảo lặng lẽ quan sát cuộc đối thoại.

Nàng tóc trắng theo bản năng, liếc nhìn cô bé. Nhưng thay vì cảm nhận được những luồng suy nghĩ phức tạp hay sự hiếu động như những học sinh khác, Inari lại thấy ở Naomi một sự tĩnh lặng đến lạ kỳ. Nó giống như một mặt hồ không một gợn sóng vào đêm trăng, một sự trống rỗng nhưng lại tràn đầy sự bao dung.

``Naomi-chan! Tới đúng lúc lắm!'' Hanabi reo lên khi thấy cô bé, cô nàng lập tức cúi xuống và bế bổng cô bé lên một cách tự nhiên. ``Naomi-chan, em xem, lớp trưởng của chúng ta lại đang `rút vỏ' nữa rồi. Em có cách nào giúp chị ấy tự tin hơn không?''

Naomi không hề tỏ ra hoảng hốt khi bị bế lên đột ngột. Cô bé nhìn thẳng vào mắt Inari, rồi chậm rãi vươn bàn tay nhỏ nhắn ra, chạm nhẹ vào mu bàn tay đang run rẩy của nàng lớp trưởng.

``Chị không cần phải cố gắng để trở nên rực rỡ như chị ấy đâu,'' em cất giọng trong trẻo, ngón tay chỉ về phía Hanabi. ``Chị giống như mặt nước phẳng lặng vậy\dots\ nhờ có sự tĩnh lặng đó mà chị mới nhìn thấu được tâm tư của mọi người. Điều đó cũng rất đáng quý mà.''

Inari sững người. Những lời nói đơn giản của cô bé tám tuổi như chạm đúng vào nút thắt trong lòng cô. Thay vì cố gắng trở thành một ``Hanabi thứ hai'', có lẽ cô nên chấp nhận và sử dụng chính sự nhạy cảm bẩm sinh của mình.

Nàng lớp trưởng khẽ hít một hơi thật sâu, đôi mắt bắt đầu ánh lên sự kiên định. Cô mỉm cười nhẹ nhàng với Naomi: ``Cảm ơn em nhé, Naomi-chan. Chị hiểu rồi.''

Hanabi cười lớn, hạ cô bé mắt xanh xuống đất và xoa đầu cô bé đầy phấn khích: ``Thấy chưa? Đến cả con bé cũng nhìn ra tài năng thiên bẩm của bà đấy! Tiến lên nào, Lớp trưởng vĩ đại!''

Naomi lẳng lặng lùi lại phía sau Kuon, nhìn Hanabi và Inari cùng nhau bước lên bục giảng. Dưới góc nhìn của em, Hanabi giống như một ngọn lửa rực rỡ luôn cháy hết mình để sưởi ấm không gian, còn Inari lại là một ngôi đền yên bình ẩn sau làn sương mù. Một người là động lực, một người là điểm tựa.

``\textit{Sự hỗ trợ lẫn nhau của họ\dots\ cũng giống như Onii-chan và Kaigou-onee-chan vậy},'' Naomi thầm nghĩ, lòng em dâng lên một cảm giác ấm áp khó tả. Ở thế giới này, sự cô đơn dường như không có chỗ đứng, vì ai cũng có một ngọn lửa của riêng mình để tựa vào.

\newpage

\subsubsection*{\phantomsection\label{ep-ER9}}
\addcontentsline{toc}{subsubsection}{{\sen{-Histoire 9-}} {\caviar Mầm xanh và Ánh nắng}}
\stpt{-Histoire 9-}\stcap{Mầm xanh và Ánh nắng}
\begin{center}{\caviar Nhân vật: \emp{lamy} \emp{aki} \emp{nene}}\end{center}

Bên khung cửa sổ tràn ngập nắng mai của lớp 1-C, Saki đang tỉ mỉ dùng chiếc kéo nhỏ tỉa lại những cành lá khô cho chậu hoa cẩm tú cầu của lớp. Với tư cách là thành viên ban chăm sóc hoa và là con gái của một người bán hoa lâu năm, cô luôn dành cho cây cỏ một sự dịu dàng đặc biệt.

``Nào Suzune-chan, đừng tưới nhiều nước quá, rễ cây sẽ bị úng đấy!'' Saki nhẹ nhàng nhắc nhở cô bạn đang cầm bình tưới với vẻ mặt đầy quyết tâm.

Suzune ngẩng lên cười hì hì: ``Tớ chỉ muốn hoa mau lớn để chúng ta có thể ngắm nhìn chúng nở rộ thôi mà! Cậu nhìn xem Saki-chan, bông này trông giống hệt như một cái bánh bao ngọt vậy!''

Đứng cạnh đó, Touka đang cầm một cuốn sổ tay nhỏ, chăm chú ghi chép lại các thông số về độ ẩm và ánh sáng. Dù là học sinh đứng đầu lớp với vẻ ngoài lạnh lùng, cô nàng tóc xanh lại không thể giấu nổi sự tò mò đối với mọi thứ xung quanh, đặc biệt là cách cô bạn thanh mai trúc mã của mình tương tác với thế giới.

Đúng lúc đó, Naomi đang đi dạo ngang qua góc lớp. Cô bé dừng lại, tò mò nhìn chậu hoa đang chớm nụ. Naomi đưa bàn tay nhỏ nhắn ra, khẽ chạm nhẹ vào lớp đất ẩm.

``Chậu cây này\dots\ đang thở rất đều,'' Naomi thì thầm, đôi mắt lục bảo ánh lên sự đồng điệu với thiên nhiên. ``Nó không có sự oán hận hay u uất như những loài cỏ dại trong rừng sâu.''

Suzune thấy cô bé thì lập tức lao tới với sự nhiệt tình thường ngày: ``Oaa! Naomi-chan! Nhìn này, chị đang trồng hoa đấy! Em có muốn tưới nước cùng tớ không? Đảm bảo vui lắm!''

Saki vội vàng can ngăn sự vồ vập của cô nàng ``Linh vật'', rồi mỉm cười dịu dàng với cô bé mắt xanh. Cô hái một bông hoa nhỏ vừa chớm nở ở một chậu khác, khẽ cài lên mái tóc tết của Naomi. ``Em nói đúng đấy, cây cỏ cũng biết cảm nhận tình cảm của con người mà.''

Touka hướng về cô bé, quan sát Naomi với ánh mắt đầy phân tích: ``Em có vẻ rất nhạy cảm với sức sống của thực vật nhỉ?''

Cô bé tóc tết ngước nhìn ba nữ sinh. Em nhận thấy ở Saki một sự chu đáo ấm áp như đất mẹ, ở Touka sự điềm tĩnh mát lành của tri thức, và ở Suzune nguồn năng lượng rực rỡ như nắng hạ.

``\textit{Chị ấy (Saki) có bàn tay của người gieo mầm sự sống,}'' Naomi thầm nghĩ khi nhận lấy bông hoa từ tay của cô gái chăm hoa. ``\textit{Những bông hoa này thật hạnh phúc khi được chị chăm sóc. Và chúng cũng thật may mắn khi có ánh nắng rực rỡ (Suzune) luôn trò chuyện cùng mỗi ngày.}''

Naomi mỉm cười rạng rỡ với bộ ba, cảm thấy rằng ngay cả những điều nhỏ bé như một chậu cây cũng có thể trở nên diệu kỳ khi được bao bọc bởi tình bạn và sự quan tâm chân thành. Thế giới người sống quả nhiên không chỉ có bê tông và đèn điện, mà còn có những mầm xanh đang vươn mình mạnh mẽ dưới ánh nắng của lòng nhân hậu.

\newpage

\subsubsection*{\phantomsection\label{ep-ER10}}
\addcontentsline{toc}{subsubsection}{{\sen{-Histoire 10-}} {\caviar Người bảo hộ và Ánh bình minh}}
\stpt{-Histoire 10-}\stcap{Người bảo hộ và Ánh bình minh}
\begin{center}{\caviar Nhân vật: \emp{miko} \emp{sora}}\end{center}

Gần khu vực cổng đền Aogami, nơi có thể nhìn bao quát cả thị trấn và cánh rừng già, Sakura và Shino đang cùng nhau tản bộ. Sakura, với trang phục nữ tu cách tân và vẻ mặt luôn tràn đầy sự hiếu kỳ, đang hào hứng chỉ cho Shino xem những mầm non đang vươn lên mạnh mẽ từ lớp đất mới.

``Cậu xem này Shino-chan, hạt giống mà tôi nghiên cứu đang phát triển rất tốt!'' Sakura reo lên, đôi mắt lấp lánh sự tự hào của một người bảo hộ khu rừng. ``Chỉ một thời gian ngắn nữa thôi, cánh rừng này sẽ rực rỡ hơn bao giờ hết.''

Shino mỉm cười rạng rỡ, nụ cười của cô mang theo một sự ấm áp thuần khiết có thể trấn an bất kỳ ai. ``Mình luôn tin tưởng vào cậu mà, Sakura-san. Cậu đã nỗ lực rất nhiều để hồi sinh nơi này.''

Đúng lúc đó, Naomi được Kaigou dẫn tới để tham quan ngôi đền. Vừa nhìn thấy Sakura, Naomi khựng lại. Cô bé cảm nhận được một luồng năng lượng quen thuộc nhưng cũng vô cùng mạnh mẽ tỏa ra từ phía cô nàng tóc hồng. Đó không phải là sát khí, mà là hơi thở của đại ngàn, là sức sống của chính khu rừng nơi Naomi từng trú ngụ trong cô độc suốt ngàn năm.

Naomi chậm rãi buông tay Kaigou, bước tới trước mặt hai người. Cô bé cúi đầu thật thấp, một cử chỉ tôn kính từ tận đáy lòng.

``Người bảo hộ\dots\ con cảm ơn Người đã che chở cho khu rừng suốt thời gian qua,'' Naomi thầm thì, giọng nói mang theo sự biết ơn sâu sắc.

Sakura đứng hình một giây, rồi bật cười tinh nghịch. Cô nàng bước tới xoa mạnh mái tóc tết của em: ``Này, đừng có gọi trịnh trọng thế chứ! Ở trường thì cứ gọi là Sakura thôi. Giờ em đã có thể tự do chạy nhảy dưới ánh mặt trời rồi, không cần phải sợ bóng tối nữa đâu nhé!''

Shino cũng quỳ một chân xuống để ngang tầm mắt với cô bé. Cô nắm lấy bàn tay nhỏ bé của Naomi, ánh mắt tràn đầy sự thấu hiểu của một người cũng từng bước ra từ vòng lặp ký ức để tìm thấy thực tại mới.

``Chào mừng em đến với thế giới này, Naomi-chan,'' Shino nói, giọng cô dịu dàng như gió xuân. ``Nơi đây sẽ luôn có chỗ cho những nụ cười chân thật.''

Naomi nhìn hai nữ sinh trước mặt. Trong mắt cô bé, Sakura giống như một cây cổ thụ vững chãi, hiên ngang bảo vệ vạn vật, còn Shino lại giống như bầu trời cao rộng, chứa đựng ánh bình minh rạng rỡ của hy vọng.

``\textit{Sakura-san mang theo linh hồn của rừng xanh, còn Shino-san mang theo ánh sáng của sự sống mới},'' cô bé thầm nghĩ, cảm nhận được một sự bình yên tuyệt đối. ``\textit{Chừng nào họ còn ở đây, thế giới này sẽ mãi luôn ổn định và tươi đẹp.}''

Cô bé tám tuổi ngẩng mặt nhìn bầu trời rộng lớn phía trên ngôi đền, hít một hơi thật sâu không khí trong lành. Hóa ra, phép màu không chỉ nằm ở những trận pháp hay bùa chú, mà nó nằm ở chính sự hiện diện của những tâm hồn cao cả luôn thầm lặng bảo vệ và sưởi ấm thế giới này mỗi ngày.

\newpage

\subsubsection*{\phantomsection\label{ep-ER11}}
\addcontentsline{toc}{subsubsection}{{\sen{-Histoire 11-}} {\caviar Những mảnh ghép của nghệ thuật}}
\stpt{-Histoire 11-}\stcap{Những mảnh ghép của nghệ thuật}
\begin{center}{\caviar Nhân vật: \emp{azki} \emp{pekora} \emp{marine}}\end{center}

Trong phòng nghệ thuật yên tĩnh của trường trung học Aogami, một bầu không khí sáng tạo đang bao trùm lấy bộ ba Azuki, Uzuki và Mari. Mỗi người đều đang đắm chìm vào thế giới nghệ thuật của riêng mình, tạo nên một sự giao thoa đầy thú vị.

Mari đang ngồi trước giá vẽ, đôi tay điêu luyện đưa những nét cọ mềm mại trên bảng màu. Cô đang hoàn thiện một bức tranh phong cảnh về thị trấn Aogami lúc hoàng hôn với những mảng màu rực rỡ và quyến rũ. Gương mặt cô hiện lên vẻ điềm tĩnh và tập trung tuyệt đối, một hình ảnh khác hẳn với sự thân thiện thường ngày.

Ở góc phòng, Uzuki đang nâng chiếc máy ảnh lên, tìm kiếm góc độ hoàn hảo để ghi lại khoảnh khắc cô nàng tóc nâu-đỏ đang mải mê với hội họa. Cô thường lang thang khắp nơi với chiếc máy ảnh, và hôm nay, cô đã tìm thấy nguồn cảm hứng trong chính những người bạn của mình.

Trong khi đó, Azuki đang đeo tai nghe, khẽ lẩm nhẩm theo một giai điệu mới. Là ca sĩ chính của câu lạc bộ âm nhạc, cô luôn tìm kiếm những âm thanh có thể truyền cảm hứng và kết nối trái tim mọi người.

Đúng lúc này, Naomi cùng Suzuki tò mò bước vào phòng nghệ thuật. Cô bé bị thu hút bởi bức tranh rực rỡ của Mari và chậm rãi tiến lại gần.

Cô nàng họa sĩ cảm nhận được sự hiện diện của em, cô quay lại mỉm cười dịu dàng: ``Ồ, Naomi-chan. Em thấy bức tranh này thế nào? Nó có giống với thị trấn mà em đang thấy không?''

Naomi nhìn chăm chú vào bức tranh, rồi khẽ chạm nhẹ vào mép giá vẽ: ``Đẹp lắm ạ\dots\ Nó giống như chị đang gom hết tất cả nắng chiều vào trong một tấm vải vậy.''

Uzuki lập tức chớp lấy cơ hội, tiếng ``tách'' của máy ảnh vang lên, ghi lại khoảnh khắc cô bé mắt xanh đang ngắm nhìn tác phẩm nghệ thuật. ``Một tấm ảnh tuyệt vời! Ánh sáng trên gương mặt Naomi-chan lúc này thật sự rất đặc biệt,'' cô hào hứng thốt lên.

Azuki tháo tai nghe ra, bước tới gần và khẽ hát một đoạn ngắn, một giai điệu trong trẻo mang theo sức mạnh của sự giải thoát và hy vọng. Naomi lắng nghe một cách say sưa, cảm thấy như tâm hồn mình đang được nâng cánh bởi âm thanh ấy.

``Mari-san lưu giữ thế giới qua màu sắc, Uzuki-san lưu giữ khoảnh khắc qua ánh sáng, còn Azuki-san lưu giữ cảm xúc qua âm thanh,'' em chậm rãi nói, đôi mắt lục bảo lấp lánh sự ngưỡng mộ.

Tới cuối buổi, Mari tặng cho cô bé một bản phác họa nhanh gương mặt của em, trong khi Uzuki hứa sẽ rửa tấm ảnh vừa rồi để tặng cô bé làm kỷ niệm.

Góc nhìn của Naomi về bộ ba này thật đặc biệt. Em cảm thấy nghệ thuật chính là cách mà con người dùng để chống lại sự lãng quên của thời gian. Mari mang đến vẻ đẹp của sự vĩnh cửu, Uzuki mang lại sự chân thật của thực tại, và Azuki mang đến sự tự do vô hình của tâm hồn.

``\textit{Thế giới này không chỉ đẹp bởi những gì chúng ta thấy, mà còn bởi cách mọi người lưu giữ nó nữa,}'' Naomi thầm nghĩ, ôm lấy bức phác họa vào lòng. Đối với cô bé từng bị mắc kẹt trong bóng tối quá khứ, nghệ thuật chính là những mảnh ghép rực rỡ giúp em cảm nhận trọn vẹn hơn sự kỳ diệu của cuộc sống người sống.

\newpage

\subsubsection*{\phantomsection\label{ep-ER12}}
\addcontentsline{toc}{subsubsection}{{\sen{-Histoire 12-}} {\caviar Trang sách và Kính hiển vi}}
\stpt{-Histoire 12-}\stcap{Trang sách và Kính hiển vi}
\begin{center}{\caviar Nhân vật: \emp{choco} \emp{haato}}\end{center}

Trong phòng chuẩn bị thí nghiệm Sinh học, Yuko đang chăm chú quan sát một mẫu vật tế bào dưới kính hiển vi. Ánh đèn phòng thí nghiệm phản chiếu trên cặp kính của cô, tạo nên một vẻ đẹp trí tuệ và có chút bí ẩn.

``Em làm xong phần phân loại tiêu bản này chưa, Akaba-san?'' cô cất giọng dịu dàng nhưng vẫn mang đầy uy lực của một giáo viên chủ nhiệm.

Koishin đang hì hục sắp xếp lại các hộp tiêu bản với vẻ mặt đầy cau có. ``Rồi ạ! Em làm cái này không phải vì muốn giúp cô đâu nhé, chẳng qua là em thấy chúng bừa bãi quá nên ngứa mắt thôi!''

Yuko khẽ mỉm cười, cô đã quá quen với tính cách ``khẩu xà tâm phật'' của cô học trò này. ``Vậy sao? Cô lại cứ tưởng em quan tâm đến tiết Sinh học chiều nay nên mới ở lại giúp cơ đấy.''

Đúng lúc đó, Naomi được Kuon dẫn vào phòng để gửi sổ điểm danh. Cô bé tò mò tiến lại gần chiếc bàn thí nghiệm, đôi mắt lục bảo nhìn chằm chằm vào những mẫu vật khô được đặt trong hộp kính.

``Đây là\dots\ những mảnh vụn của thời gian sao ạ?'' Naomi khẽ hỏi, ngón tay nhỏ nhắn chỉ vào một mẫu xác ve sầu khô.

Koishin giật nảy mình khi thấy cô bé, cô nàng lập tức đỏ mặt, che đi hộp tiêu bản mình vừa mới cẩn thận lau chùi. ``Này nhóc! Vào phòng thí nghiệm phải gõ cửa chứ! Với lại đây không phải là mảnh vụn thời gian gì hết, nó là mẫu vật sinh học, hiểu chưa?''

Nói đoạn, cô lục lọi trong túi áo, lôi ra một chiếc bánh quy kẹp kem nhỏ và nhét vào tay Naomi. ``Cầm lấy đi! Tôi thấy nó vướng víu quá nên mới cho em đấy, đừng có mà hiểu lầm là tôi thích trẻ con!''

Yuko-sensei đứng cạnh quan sát, ánh mắt cô hiện lên sự ấm áp. Cô cúi xuống, điều chỉnh lại tiêu cự kính hiển vi rồi mời Naomi nhìn vào. ``Nhìn xem Naomi-chan, sự sống đôi khi ẩn chứa trong những điều nhỏ bé nhất mà mắt thường không thể thấy được.''

Cô bé nhìn vào kính hiển vi, rồi ngước lên nhìn hai người. Cô bé nhận thấy ở Yuko-sensei một sự điềm tĩnh và thấu hiểu sâu sắc về quy luật của thế giới tự nhiên, còn ở Koishin là một sự mâu thuẫn đầy thú vị của con người thời hiện đại.

``\textit{Cô giáo Yuko giống như một nhà khoa học đang giải mã bí mật của sự sống,}'' Naomi thầm nghĩ khi nhấm nháp miếng bánh quy ngọt lịm. ``\textit{Còn Koishin-san lại giống như một cuốn sách có bìa cứng và gai góc, nhưng bên trong lại chứa đầy những câu chuyện dịu dàng.}''

Naomi mỉm cười cảm ơn rồi theo Suzuki rời khỏi phòng. Cô bé cảm thấy rằng, dù là dưới ống kính hiển vi hay sau những lời nói sắc sảo, lòng tốt và sự chân thành vẫn luôn là thứ ngôn ngữ dễ hiểu nhất, giúp em xua tan đi cảm giác lạc lõng trong thế giới mới mẻ này.

\newpage

\subsubsection*{\phantomsection\label{ep-ER13}}
\addcontentsline{toc}{subsubsection}{{\sen{-Histoire 13-}} {\caviar Trò đùa sau cánh cửa}}
\stpt{-Histoire 13-}\stcap{Trò đùa sau cánh cửa}
\begin{center}{\caviar Nhân vật: \emp{suisei} \emp{noel}}\end{center}

\begin{center}
  (Lấy cảm hứng từ \href{https://x.com/suisei_hosimati/status/1480161260058845187/video/1}{cảnh Behind the Scene của \textit{hololive ERROR} giữa Suisei và Noel được đăng trên Twitter/X.})
  \pov{-}
\end{center}

Ánh hoàng hôn nhuộm đỏ hành lang trường trung học Aogami, báo hiệu một ngày học nữa đã kết thúc. Trong lớp 1-C, Nanase đang thong thả thu xếp sách vở vào cặp. Ở bên cạnh, Suzuki đang dắt tay Naomi, cả hai đang đợi cô nàng tóc xanh để cùng ra về.

Đúng lúc này, từ phía ngoài cửa lớp vang lên tiếng đập cửa rầm rầm.

``\textbf{Nanase-chan! Mở cửa cho tớ với!!}'' Yae đang đứng bên ngoài, gương mặt dán sát vào lớp kính cửa, hai tay không ngừng đập vào cánh cửa trượt đã bị khóa chặt.

\begin{figure}[H]
  \centering
  \includegraphics[width=12.5cm]{../../../../Image/Book?/eb.png}
  \caption{(Cảnh giống như vậy, nhưng đặt ở bối cảnh thời hiện đại.)}
\end{figure}

Nanase nhướng mày, khẽ nhếch mép cười đầy tinh quái. Cô đứng tựa lưng vào cửa, thản nhiên hỏi: ``Sao vậy, Yae-san? Có chuyện gì mà cậu có vẻ vội vàng thế?''

``\textbf{Nè! Cái cửa bị khóa rồi! Thế này là thế nào đây!? Mở cửa cho tớ vào với!}'' Yae gào lên, giọng nói đầy vẻ ấm ức.

Cô bạn khoanh tay trước ngực, giọng điệu vô cùng bình thản: ``Sao, có chuyện gì?''

``Tớ để quên cặp trong đó nên không về được! Nanase-chan làm ơn đi mà!'' cô nàng Thể chất nhảy tưng tưng ở bên ngoài, điệu bộ luống cuống của cô nàng trông giống hệt như một chú gấu lớn đang bị nhốt ngoài hang.

Nanase không nhịn được cười, cô liếc nhìn chiếc cặp của Yae vẫn còn nằm chơ vơ trên bàn. ``Tôi không cho cậu vào đâu. Cậu cứ thế mà về đi, cho chừa cái tật quên đồ. Một cô gái lý tưởng như cậu nói mà lại hậu đậu thế sao?''

``\textbf{Sao kỳ zậy!?? Nanase-chan là đồ ác độc!}'' Yae tiếp tục múa may quay cuồng bên ngoài cửa, tạo nên một cảnh tượng vô cùng hài hước.

Chứng kiến cảnh tượng đó, cả Nanase và Suzuki đều bật cười thành tiếng, cô bạn tóc xanh thậm chí còn bắt chước điệu nhảy cà tưng của cô nàng Thể dục tội nghiệp. Ngay cả Naomi, vốn thường ngày rất trầm mặc, lúc này cũng không khỏi tò mò dõi theo từng chuyển động của Yae qua lớp kính.

Cô bé nghiêng đầu, đôi mắt lục bảo quan sát sự tương tác giữa hai người bạn thân. ``\textit{Yae-san giống như một chú gấu lớn đang nhảy múa vì mất đi kho báu của mình,}'' Naomi thầm nghĩ, một tia sáng vui vẻ lướt qua mắt em. ``\textit{Còn Nanase-san lại giống như một người quản trò đang cười thầm sau tấm màn.}''

Naomi bước lại gần cửa, vẫy tay nhẹ với Yae đang mếu máo ở bên ngoài. Cô bé cảm nhận được một luồng năng lượng rất rộn ràng và ấm áp tỏa ra từ hai người. Đó không phải là sự tranh cãi gay gắt, mà là một sự gắn kết sâu sắc đến mức họ có thể thoải mái trêu đùa nhau như vậy.

Cuối cùng, Nanase cũng chịu mở khóa cửa. Yae lao vút vào trong như một cơn lốc, vơ lấy chiếc cặp rồi quay sang lườm cô bạn thơ ấu một cái sắc lẹm, nhưng ngay sau đó lại cười xòa khi thấy Naomi đang nhìn mình.

``\textit{Trêu chọc nhau cũng là một cách để bày tỏ tình yêu thương sao?}'' Naomi tự hỏi, lòng em cảm thấy nhẹ nhàng lạ kỳ. Hóa ra, ở thế giới này, ngay cả những trò đùa oái oăm nhất cũng có thể chứa đựng những mảnh ghép ngọt ngào của tình bạn.

\newpage

\subsubsection*{\phantomsection\label{ep-ER14}}
\addcontentsline{toc}{subsubsection}{{\sen{-Histoire 14-}} {\caviar Những mảnh ghép của ánh sáng}}
\stpt{-Histoire 14-}\stcap{Những mảnh ghép của ánh sáng}

Con đường dẫn về nhà Hikawa hôm nay dường như ngắn hơn thường lệ dưới ánh hoàng hôn vàng vọt. Sau một ngày dài náo nhiệt tại trường trung học Aogami, bộ ba Hikawa cùng cô bé Naomi đang thong thả bước đi trên con phố nhỏ. Tiếng ve sầu kêu râm ran cuối hè tạo nên một bản nhạc nền êm ả cho buổi chiều tàn.

Suzuki, người luôn là người khơi mào các cuộc trò chuyện, khẽ cúi xuống nhìn cô bé mắt xanh đang ngoan ngoãn nắm lấy tay mình. ``\vie{Này Naomi-chan, hôm nay em ở lớp thấy thế nào? Có thấy các anh chị trong lớp thú vị không?}'' cô hào hứng hỏi, đôi mắt lấp lánh sự tò mò.

Kaigou cũng mỉm cười, giọng cô dịu dàng hòa vào cơn gió nhẹ: ``\vie{Đúng đấy, em tiếp xúc với nhiều người như vậy, có ấn tượng đặc biệt với ai nhất không?}''

Kuon im lặng bước bên cạnh, nhưng ánh mắt anh luôn dõi theo từng biểu cảm của ba người. Cậu cũng muốn biết liệu thế giới loài người có thực sự mang lại niềm vui cho cô bé sau ngàn năm cô độc hay không.

Naomi dừng lại, đôi mắt lục bảo nhìn về phía chân trời xa xăm, nơi những áng mây đang chuyển dần sang màu tím thẫm. Em suy nghĩ một chút, rồi bắt đầu kể bằng giọng nói trong trẻo của mình.

``\vie{Mọi người ở đó\dots\ thật kỳ lạ,}'' em bắt đầu nói. ``\vie{Honoka-san thì ồn ào như tiếng pháo hoa, Shiina-san thì ngầu như một chiến binh thực thụ. Sakura-san mang theo hơi thở của rừng xanh, còn Shino-san lại giống như ánh bình minh vậy\dots}''

Cô bé kể lại từng mảnh ký ức nhỏ trong ngày: từ mẩu sô-cô-la của Miku, hộp sữa của Kaoru, cho đến những trò đùa nghịch ngợm của Nanase và Yae. Mỗi cái tên, mỗi gương mặt đều được em ghi nhớ một cách trân trọng.

``\vie{Lúc đầu, em đã rất sợ rằng mình sẽ không thuộc về nơi này,}'' Naomi siết chặt lấy tay Suzuki. ``\vie{Nhưng hóa ra, mỗi người trong lớp đều là một mảnh ghép của ánh sáng. Họ có những màu sắc khác nhau, nhưng đều mang lại sự ấm áp. Khi ở cạnh họ, em không còn thấy lạnh lẽo nữa.}''

Nghe những lời tâm sự chân thành đó, Kuon khẽ xoa đầu cô bé. ``\vie{Anh đã bảo mà. Mọi người đều yêu quý em vì em chính là Naomi.}''

Kaigou mỉm cười đồng tình: ``\vie{Đúng thế. Ngày mai và cả những ngày sau nữa, em sẽ còn khám phá ra thêm nhiều điều tuyệt vời hơn ở ngôi trường này đấy.}''

Suzuki cười toe toét, bế bổng Naomi lên vai như thường lệ: ``\vie{Vậy thì từ mai em lại đến lớp nhé? Lớp 1-C mà thiếu em là buồn lắm đấy!}''

Cô bé gật đầu, nở một nụ cười rạng rỡ nhất từ trước đến nay. Nụ cười ấy xua tan đi những tàn dư cuối cùng của nỗi u uất ngàn năm.

``\vie{Vâng, ngày mai em lại muốn đến trường ạ.}''

Thị trấn Aogami có thể vẫn còn đó những bí ẩn chưa được giải mã, những nghi thức cổ xưa đang chờ đợi ở mùa Lễ hội sắp tới. Nhưng ngay tại lúc này, trong căn phòng nhỏ này, niềm hạnh phúc giản đơn nhất đang nảy nở. Một sự bình yên thực sự sau cơn bão lớn.

Hành trình của chúng tôi, với sự góp mặt của thành viên nhỏ tuổi nhất - Naomi, có lẽ sẽ còn nhiều gian nan phía trước. Nhưng chừng nào chúng tôi còn ở bên nhau, cùng chia sẻ những lo âu và nụ cười như thế này, tôi tin rằng mọi thử thách đều có thể vượt qua.

Bốn bóng người đổ dài trên mặt đường, cùng nhau hướng về phía ngôi nhà ngập tràn tình thương. Nhưng mảnh truyện nhỏ về những ngày đầu tiên của Naomi tại Aogami khép lại, nhưng hành trình hòa nhập và tìm thấy hạnh phúc thực sự của cô bé mắt xanh chỉ mới vừa bắt đầu.

\end{document}